% SPDX-License-Identifier: CC-BY-SA-4.0
% Part: Introduction
% Section: Words and languages
% Exercise: Palindromes and Fibonacci sequence

\begingroup

\begin{exercice}[Palindrome]\label{exo:introduction/languages/palindrome}
  Un palindrome sur un alphabet $\Sigma$ est un mot pouvant être lu indifféremment de gauche à droite ou de droite à gauche (par exemple : Laval, été).
  Formellement, un mot $u \in \Sigma^\star$ est un palindrome si
  $$\forall i \in \llbracket |u| \rrbracket, u[i] = u[1+|u| - i].$$

  On définit sur $\Sigma =\{a,b\}$ la suite de mots $(f_n)_{n>0}$ (suite de Fibonacci) de la façon suivante :
  $$
  \left\{\begin{array}{rcll}
  f_1 &=& a\\
  f_2 &=& ab \\
  f_{n+2} &=& f_{n+1}\cdot f_n & \text{si $n>0$}
  \end{array}
  \right.
  $$
  
  \begin{question}
  \item Montrer que pour tout $i \ge 2$, il existe un palindrome $u_i$ tel que:
    $$ f_i = \left \{\begin{array}{ll}
    u_i\cdot a\cdot b & \mbox{si $i$ pair} \\
    u_i\cdot b\cdot a & \mbox{si $i$ impair}
  \end{array}\right.$$
  \end{question}

\end{exercice}

\begin{correction}

  \begin{question}
  \item Preuve par récurrence forte. Soit $i \ge 2$. Supposons la propriété pour tout $k<i$,
    et prouvons la propriété pour $i$. On sépare les cas $i\le 5$ et $i > 5$ pair et $i > 5$ impair.
    \begin{description}
    \item [Cas $i \le 5$.]  On a $f_2 = ab$ ($u_2=\varepsilon$), $f_3 = aba$ ($u_3=a$) , $f_4 = abaab$ ($u_4=aba$) car $\varepsilon$.
    \item [Cas $i > 5$ impair.] $n-2$ est impair et $n-3$ est pair et
      par hypothèse de récurrence, $f_{n-2} = u_{n-2} b a$ et $f_{n-3}= u_{n-3} a b$
      avec $u_{n-2}$ et $u_{n-3}$ des palindromes.

      Donc $f_{n} = f_{n-1} f_{n-2} = f_{n-2} f_{n-3} f_{n-2} = u_{n-2} b a u_{n-3} a b u_{n-2} b a$.

      Mais $u_{n-2}$ et $u_{n-3}$ sont des palindromes, donc $u_{n-2} b a u_{n-3} a b u_{n-2}$ est également un palindrome.

      En posant $u_{n} = u_{n-2} b a u_{n-3} a b u_{n-2}$, on a alors $f_{n}  = u_n ba$ où $u_n$ est un palindrome.

    \item [Cas $i > 5$ pair.]
      $n-2$ est pair et $n-3$ est impair et
      par hypothèse de récurrence, $f_{n-2} = u_{n-2} a b$ et $f_{n-3}= u_{n-3} b a$
      avec $u_{n-2}$ et $u_{n-3}$ des palindromes.
      
      En posant $u_{n} = u_{n-2} a b  u_{n-3} b a  u_{n-2}$, on a alors $f_{n}  = u_n ab$ où $u_n$ est un palindrome.
    \end{description}
  \end{question}

\end{correction}

\endgroup
\endinput
